\documentclass[11pt]{article}
\usepackage[margin=1.5cm]{geometry}
\usepackage{amsmath}
\usepackage{amssymb}
\usepackage{enumitem}
\usepackage{array}
\usepackage{fancyvrb}
\setlength{\columnsep}{0.7cm}
\setlist{nosep,leftmargin=*}
\setlength{\parindent}{0pt}
\pagestyle{empty}

\begin{document}
\twocolumn

\title{In-Class Exercise: Recursive Binary Search}
\author{Prof. Jordan C. Hanson}
\date{\today}

\maketitle

\section{The Idea of Binary Search}
Binary search finds a target in a \textbf{sorted} list by repeatedly checking the midpoint and discarding the half that cannot contain the target. Use the list
\begin{verbatim}
values = [3, 7, 11, 18, 23,
          29, 34, 42, 47, 53,
          61, 68, 74, 83, 91]
\end{verbatim}
The indices run from \(0\) through \(14\).  Suppose the target is \(74\). Initially,
\[
\texttt{low}=0,\qquad \texttt{high}=14.
\]
Compute the midpoint index with
\[
\texttt{mid}=\left\lfloor\frac{\texttt{low}+\texttt{high}}{2}\right\rfloor.
\]

The oddly shaped brackets indicate the \textit{floor function}, or ``rounding down.''  Complete the trace:
\begin{center}
\renewcommand{\arraystretch}{1.25}
\begin{tabular}{c|c|c|c|c}
\textbf{low} & \textbf{high} & \textbf{mid} & \textbf{value} & \textbf{keep}\\ \hline
0 & 14 & \rule{0.55cm}{0.15mm} & \rule{0.7cm}{0.15mm} & \rule{1.0cm}{0.15mm}\\
\rule{0.5cm}{0.15mm} & \rule{0.5cm}{0.15mm} & \rule{0.55cm}{0.15mm} & \rule{0.7cm}{0.15mm} & \rule{1.0cm}{0.15mm}\\
\rule{0.5cm}{0.15mm} & \rule{0.5cm}{0.15mm} & \rule{0.55cm}{0.15mm} & \rule{0.7cm}{0.15mm} & \rule{1.0cm}{0.15mm}\\
\rule{0.5cm}{0.15mm} & \rule{0.5cm}{0.15mm} & \rule{0.55cm}{0.15mm} & \rule{0.7cm}{0.15mm} & \rule{1.0cm}{0.15mm}\\
\end{tabular}
\end{center}

What index should the search return? \rule{1.2cm}{0.15mm}

\section{Base Case and Recursive Cases}
A recursive binary search keeps the original list and changes only the indices \texttt{low} and \texttt{high}. The \textbf{base case} for an unsuccessful search occurs when \texttt{low} $>$ \texttt{high}. There are then no indices left to search, so return \texttt{-1}.  If \texttt{values[mid]} equals the target, return \texttt{mid}. Otherwise: \\
\begin{itemize}
    \item if the target is smaller, search \texttt{low} through \texttt{mid-1};
    \item if the target is larger, search \texttt{mid+1} through \texttt{high}.
\end{itemize}

\section{Complete the Recursive Function}
Fill in the missing expressions.

\begin{Verbatim}[fontsize=\small]
def binary_search(values,target,low,high):
    if ____________________:
        return -1

    mid = __________________

    if values[mid] == target:
        return ______________

    if target < values[mid]:
        return binary_search(values,target,low,____)

    return binary_search(values,target,____,high)
\end{Verbatim}

Call the function using the entire list:
\begin{Verbatim}[fontsize=\small]
result = binary_search(values,74,0,len(values)-1)
print(result)
\end{Verbatim}

Predict the output before running it: \rule{1.5cm}{0.15mm}

\section{Search for a Missing Value}
Now search for \(50\):
\begin{verbatim}
result = binary_search(values,50,0,len(values)-1)
\end{verbatim}

Trace the sequence of midpoint values examined.  What value should the function return?  Explain why \texttt{low > high} eventually becomes true:
\vspace{1cm}


\section{Why Binary Search Is Efficient}
Each recursive call discards about half of the remaining search region.
A list of roughly \(1024=2^{10}\) sorted values therefore needs at most
about \(10\) midpoint decisions, rather than checking every value one at
a time. \\

\textbf{Question:} Why must the list be sorted for this method to work?
\vspace{0.55cm}
\hrule

\end{document}
